\documentclass[11pt]{article}

\usepackage[margin=1in]{geometry}
\usepackage{amsmath,amssymb,amsthm}
\usepackage{booktabs}
\usepackage{hyperref}
\usepackage{algorithm}
\usepackage{algorithmic}

\newtheorem{proposition}{Proposition}
\newcommand{\R}{\mathbb{R}}
\newcommand{\Z}{\mathbb{Z}}
\newcommand{\BiCut}{\textsc{BiCut}}

\title{\BiCut: Checked Single-Level Reformulations\\for Bilevel Programs with Integer Leaders}
\author{BiCut contributors}
\date{}

\begin{document}
\maketitle

\begin{abstract}
\BiCut\ solves bilevel linear programs with an integer leader through
single-level mixed-integer programs and checks every answer it returns. For a
continuous follower it uses the big-M KKT reformulation with complementarity
on both the constraint slacks and the variable reduced costs. We show that
every feasible point of this model is follower-optimal for any value of the
big-M constant, so an incorrect constant can only make the model infeasible
or its optimum too high. On 122 generated instances with $M=1000$, HiGHS
solves 119 to optimality, every returned point passes a one-LP follower
check, and on the 34 instances small enough to enumerate the leader the KKT
optimum equals the enumerated optimum. Sweeping $M$ from 1 to $10^4$ never
produced a follower-suboptimal answer; at $M=10$, 70 of 117 runs reported
infeasibility and 7 returned a feasible but suboptimal value. For an integer
follower, the KKT model of the follower's LP relaxation returns checked
bilevel-feasible points, and an exact loop over leader decisions bounded by
the high-point relaxation and warm-started from those points proves
optimality on 35 of the 80 smallest optimally solved BOBILib instances
within 60 seconds each; all 35 proofs agree with the optimum BOBILib
reports, and 49 of its final values equal it.
\end{abstract}

\section{Introduction}
\label{sec:intro}

A bilevel program optimizes a leader's decision $x$ subject to the
condition that a follower responds optimally:
\begin{align}
\min_{x,\,y}\quad & d^\top x + e^\top y \label{eq:bilevel}\\
\text{s.t.}\quad & Cx + Dy \le h,\quad x \in X \cap \Z^{n_x}, \nonumber\\
& y \in S(x) = \operatorname*{argmin}_{y' \ge 0}\{\, c^\top y' : Ay' \le b + Bx \,\}. \nonumber
\end{align}
We consider the optimistic version, in which the leader picks the best $y$
among the follower's optimal responses. Applications include interdiction,
pricing and Stackelberg games \cite{colson2007overview, dempe2002foundations}.

The standard way to solve~\eqref{eq:bilevel} with an off-the-shelf MILP solver
replaces $y \in S(x)$ by the follower's KKT conditions and linearizes the
complementarity products with a constant $M$
\cite{fortuny1981representation}. Choosing $M$ is the weak point: bound
tightening can produce invalid values \cite{pineda2019solving}, and verifying
that a given $M$ is valid is as hard as solving the bilevel program
\cite{kleinert2020nofreelunch}. When the follower has integer variables the KKT
conditions do not characterize optimality at all, and dedicated
branch-and-cut solvers are used instead
\cite{denegre2009branch, fischetti2017new, tahernejad2020branch}.

This paper makes one claim: a single-level reformulation becomes a dependable
tool when each of its answers is checked and its failure modes are known in
advance. For an LP follower the failure mode of a bad $M$ is restricted to
infeasibility or a too-high objective, never a follower-suboptimal point; for
an integer follower a checked relaxation-based answer is a valid incumbent
for an exact leader enumeration. Contributions:
\begin{enumerate}
\item A statement of the big-M KKT model with both complementarity families
  and a proof that its feasible points are bilevel feasible for every $M>0$
  (Proposition~\ref{prop:anyM}).
\item An evaluation on 122 generated instances: correctness against
  enumeration, a one-LP answer check, and a sweep over $M$
  (Section~\ref{sec:eval-lp}).
\item For integer followers, a checked KKT heuristic and an HPR-bounded
  leader enumeration with an optimality proof, evaluated on BOBILib instances
  with known optima (Section~\ref{sec:eval-int}).
\item An open-source implementation: Python reference code on HiGHS and SCIP,
  and a Rust workspace with problem types, MPS/LP writers and a command-line
  tool.
\end{enumerate}

\section{Background and Related Work}
\label{sec:related}

\paragraph{KKT and big-M.} Replacing the follower LP by its KKT system and
encoding complementarity with binary variables goes back to Fortuny-Amat and
McCarl~\cite{fortuny1981representation}. Pineda and
Morales~\cite{pineda2019solving} show that big-M values from common bound
tightening heuristics can cut off the bilevel optimum, and Kleinert et
al.~\cite{kleinert2020nofreelunch} prove that verifying a big-M is as hard as
the original problem. BilevelJuMP~\cite{diasgarcia2024bileveljump} offers
KKT, strong-duality and SOS1/indicator encodings in Julia. We add a
measurement of how often a given $M$ fails and in which way.

\paragraph{Integer followers.} With an integer follower, exact methods
include the branch-and-bound algorithm of Moore and Bard~\cite{moore1990mixed},
branch-and-cut by DeNegre and Ralphs~\cite{denegre2009branch} and its
implementation in MibS~\cite{tahernejad2020branch}, and the branch-and-cut
algorithm of Fischetti et al.~\cite{fischetti2017new}, which separates
intersection cuts~\cite{balas1971intersection} derived from bilevel-free
sets~\cite{fischetti2016intersection, fischetti2018use}. Kleinert et
al.~\cite{kleinert2021survey} survey MILP techniques for bilevel optimization.
The leader enumeration in Section~\ref{sec:enum} is a simple reference method
for pure-integer leaders; it is not competitive with these solvers on large
instances, and we use it to obtain proven optima with standard MIP calls.

\paragraph{Benchmarks.} BOBILib~\cite{thurauf2026bobilib} collects bilevel
instances in an MPS-plus-auxiliary-file format together with the best known
or optimal objective values, which we use as reference.

\section{Approach}
\label{sec:approach}

\subsection{Big-M KKT model for an LP follower}

Let the follower rows be extended by any finite bounds $y \le u$, giving
$A'y \le b' + B'x$ with $m'$ rows. The follower LP is optimal at $y$ for a
fixed $x$ if and only if there is $\lambda \ge 0$ with reduced cost
$r = A'^\top\lambda + c \ge 0$, $\lambda_k (b' + B'x - A'y)_k = 0$ for all
$k$, and $y_j r_j = 0$ for all $j$. With binaries $z \in \{0,1\}^{m'}$ and
$t \in \{0,1\}^{n_y}$ the model is
\begin{align}
\min\quad & d^\top x + e^\top y \nonumber\\
\text{s.t.}\quad & Cx + Dy \le h,\quad A'y - B'x \le b',\quad A'^\top\lambda + c \ge 0, \nonumber\\
& \lambda \le M z,\quad b' + B'x - A'y \le M(\mathbf{1} - z), \label{eq:kkt}\\
& y \le M t,\quad A'^\top\lambda + c \le M(\mathbf{1} - t), \nonumber\\
& x \in X \cap \Z^{n_x},\ y \ge 0,\ \lambda \ge 0,\ z, t \text{ binary}. \nonumber
\end{align}

\begin{proposition}
\label{prop:anyM}
For every $M > 0$, the $(x,y)$ part of every feasible point of~\eqref{eq:kkt}
is feasible for~\eqref{eq:bilevel}. Hence the optimum of~\eqref{eq:kkt} is
either infeasible or at least the bilevel optimum, and it equals the bilevel
optimum whenever $M$ bounds $\lambda$, the slacks, $y$ and $r$ at some optimal
bilevel solution with a matching multiplier.
\end{proposition}
\begin{proof}
If $z_k = 0$ then $\lambda_k \le 0$, so $\lambda_k = 0$; if $z_k = 1$ the
slack of row $k$ is at most $0$ and, by primal feasibility, equals $0$. The
same argument with $t_j$ gives $y_j r_j = 0$. So $(y,\lambda)$ satisfies the
KKT conditions of the follower LP at $x$, which are sufficient for
optimality; the leader constraints are part of the model. The second
statement follows because the model is then a restriction of the bilevel
program, and a bounded optimal pair is feasible for it.
\end{proof}

A consequence is that a follower check cannot detect a bad $M$ for this
model: the answer is always bilevel feasible and the error, if any, is in
the leader objective. We still run the check, one LP per answer, as a guard
against numerical tolerance and modelling errors. The check accepts $(x,y)$
if all constraints hold and $c^\top y \le \varphi(x) + 10^{-6}(1+|\varphi(x)|)$,
where $\varphi(x)$ is the follower LP value at $x$.

\subsection{Integer follower: checked KKT of the LP relaxation}

When some follower variables are integer, we apply the KKT model to the
follower's LP relaxation and keep the integrality of $y$ in the MILP. This is
the KKT method in Section~\ref{sec:eval-int}; it uses indicator constraints
in SCIP instead of big-M.

\begin{proposition}
\label{prop:relax}
If $(x,y)$ satisfies the KKT conditions of the LP relaxation of the follower
at $x$ and $y$ is integral, then $y$ is optimal for the integer follower at
$x$.
\end{proposition}
\begin{proof}
$\varphi_{\mathrm{LP}}(x) \le \varphi_{\mathrm{IP}}(x) \le c^\top y =
\varphi_{\mathrm{LP}}(x)$.
\end{proof}
The model therefore returns bilevel-feasible points, but it excludes every
integer response whose value exceeds the LP value, so it can return a
suboptimal objective or report infeasibility. We check each answer by
solving the follower MIP at the returned $x$.

\subsection{HPR-bounded leader enumeration}
\label{sec:enum}

For a pure-integer leader with finite bounds we use
Algorithm~\ref{alg:enum}. The high-point relaxation (HPR) is
\eqref{eq:bilevel} without the condition $y \in S(x)$.

\begin{algorithm}[t]
\caption{HPR-bounded leader enumeration}
\label{alg:enum}
\begin{algorithmic}[1]
\REQUIRE bilevel instance, optional incumbent value $U$ of a checked solution
\STATE $U \leftarrow U$ or $+\infty$; $\mathcal{E} \leftarrow \emptyset$
\LOOP
  \STATE solve the HPR with the leader vectors in $\mathcal{E}$ excluded
  \IF{infeasible \OR its value $L \ge U$} \RETURN $U$ (optimal) \ENDIF
  \STATE $\hat{x} \leftarrow$ leader part of the HPR optimum
  \STATE $\varphi \leftarrow$ follower MIP value at $\hat{x}$
  \STATE $F \leftarrow \min\{d^\top\hat{x} + e^\top y : Cx + Dy \le h,\ Ay \le b + B\hat{x},\ c^\top y \le \varphi,\ y \text{ follower-feasible}\}$
  \STATE $U \leftarrow \min(U, F)$; $\mathcal{E} \leftarrow \mathcal{E} \cup \{\hat{x}\}$
\ENDLOOP
\end{algorithmic}
\end{algorithm}

A leader vector $\hat{x}$ is excluded with one binary per direction and bound,
$x_i \le \hat{x}_i - 1$ or $x_i \ge \hat{x}_i + 1$ for at least one $i$, which
is valid for general integers.

\begin{proposition}
If every sub-MIP is solved to optimality, Algorithm~\ref{alg:enum}
terminates and returns the optimistic bilevel optimum.
\end{proposition}
\begin{proof}
$F(\hat{x})$ is the optimistic value at $\hat{x}$, so $U$ is always the value
of a bilevel-feasible point (or the given incumbent). The bilevel optimum is
at least $\min(U, L)$, since every leader vector outside $\mathcal{E}$ has
optimistic value at least the HPR value over that set. The loop stops only
when $L \ge U$ or the set is empty. Each iteration removes one of finitely
many leader vectors.
\end{proof}
The implementation reports a proof only when the loop stops in this way and
no sub-MIP hit its time limit.

\section{Implementation}
\label{sec:impl}

The reference implementation is in Python. \texttt{kkt\_highs.py} builds
model~\eqref{eq:kkt} in HiGHS~\cite{huangfu2018highs};
\texttt{kkt\_eval.py} runs the LP-follower evaluation;
\texttt{bobilib\_eval.py} reads BOBILib MPS and auxiliary files and runs the
integer-follower methods in SCIP~\cite{bolusani2024scip}. The repository also
contains a Rust workspace with typed bilevel problem structures, MPS and LP
readers and writers, reformulation passes as library code with 718 unit
tests, and a command-line tool that reads a bilevel problem in TOML and
writes its high-point relaxation as MPS or LP.

\section{Evaluation}
\label{sec:eval}

We ask three questions.
\textbf{RQ1}: does the KKT model with a fixed $M$ return the bilevel optimum
on instances with an LP follower?
\textbf{RQ2}: how does it fail when $M$ is too small?
\textbf{RQ3}: on integer-follower instances with known optima, how often do
the checked KKT heuristic and the leader enumeration reach the optimum?

\paragraph{Setup.} HiGHS 1.15.1 and SCIP 10.0 (PySCIPOpt 6.2.1), Python 3.14,
on an Apple M2 shared with other jobs; wall-clock times are therefore
indicative. KKT solves use a 60\,s limit for RQ1 and 30\,s for RQ2; each
BOBILib method has 60\,s per instance.

\subsection{RQ1 and RQ2: LP follower}
\label{sec:eval-lp}

\paragraph{Instances.} 122 instances in four generated families
(\texttt{generators.py}): 72 knapsack-type instances in which a binary leader
removes capacity from an LP knapsack with weakly correlated, strongly
correlated or uncorrelated weights ($n = 8$ to $50$); 28 dense random
instances with integer $x \in [0,3]^{n_x}$ ($5 \le n_x \le 20$,
$5 \le n_y \le 30$); 12 instances with pairwise follower coupling
($n = 5$ to $25$); and 10 target-defence games in which a binary defender
lowers attack bounds and an LP attacker maximizes damage ($n = 5$ to $20$).
For the 34 instances with at most 4096 leader vectors we compute the exact
optimum by solving one LP per leader vector.

\begin{table}[t]
\centering
\caption{KKT model with $M = 1000$ on 122 generated LP-follower instances.
``Checked'': the returned point passes the follower-LP check. ``Exact'':
instances whose leader was enumerated; ``$=$'': KKT optimum equals the
enumerated optimum. Nodes and times are means over solved instances.}
\label{tab:rq1}
\begin{tabular}{lrrrrrrr}
\toprule
Family & $n$ & Optimal & Checked & Exact & $=$ & Nodes & Time (s) \\
\midrule
Knapsack-type      & 72 & 72 & 72 & 18 & 18 & 1.0 & 0.25 \\
Dense random       & 28 & 25 & 25 &  4 &  4 & 438.9 & 8.55 \\
Pairwise coupling  & 12 & 12 & 12 &  6 &  6 & 0.0 & 0.01 \\
Target defence     & 10 & 10 & 10 &  6 &  6 & 1.0 & 0.19 \\
\midrule
Total & 122 & 119 & 119 & 34 & 34 & & \\
\bottomrule
\end{tabular}
\end{table}

\paragraph{RQ1.} Table~\ref{tab:rq1} shows that with $M = 1000$ HiGHS solves
119 of the 122 instances to optimality, every returned point passes the
follower check, and the KKT optimum equals the enumerated optimum on all 34
enumerated instances. Of the three dense instances that are not solved, one
is reported infeasible for every $M$ in the sweep and two reach the time
limit with incumbents that pass the check. The linear relaxation of the KKT
model is weak outside the knapsack family: its mean gap to the optimum is
47\% on the dense family and larger on the coupling and defence families,
while the knapsack family has zero gap because leader and follower share the
objective there and the high-point relaxation already attains the bilevel
optimum. The dense family accordingly needs branching (mean 439 nodes).

\begin{table}[t]
\centering
\caption{Outcome of the KKT model for different $M$ on the 117 instances
with a reference optimum (the enumerated optimum, or the checked optimum at
$M = 10^4$).}
\label{tab:rq2}
\begin{tabular}{rrrrrr}
\toprule
$M$ & Correct & Infeasible & Suboptimal & Follower-suboptimal & Time limit \\
\midrule
1, 2, 5  &   0 & 117 & 0 & 0 & 0 \\
10       &  40 &  70 & 7 & 0 & 0 \\
50       & 115 &   1 & 0 & 0 & 1 \\
100      & 116 &   0 & 0 & 0 & 1 \\
1000, $10^4$ & 117 & 0 & 0 & 0 & 0 \\
\bottomrule
\end{tabular}
\end{table}

\paragraph{RQ2.} Table~\ref{tab:rq2} confirms Proposition~\ref{prop:anyM} in
practice: of the 976 sweep runs over all 122 instances, 527 returned a
point and all 527 pass the follower check. A too-small $M$
shows up in two ways. Most often the model becomes infeasible (every run at
$M \le 5$, 70 runs at $M = 10$, one knapsack instance at $M = 50$). Less
often it returns a bilevel-feasible point with a worse leader value: seven
dense instances at $M = 10$, with objective errors between 14\% and 83\% of
the optimum. The second case is the dangerous one, because the answer passes
every feasibility check; only a run with a larger $M$ or a lower bound
reveals it. For these instances every $M \ge 1000$ was sufficient.

\subsection{RQ3: integer follower on BOBILib}
\label{sec:eval-int}

\paragraph{Instances.} We take the 80 smallest general-bilevel instances,
by compressed MPS size, that BOBILib reports as solved to optimality: 50 from
the DeNegre \texttt{miblp} set, 10 Xu--Wang \texttt{bmilplib} instances with
a mixed-integer follower, 11 MIPLIB-derived instances (\texttt{p0033},
\texttt{enigma}, \texttt{stein27}, \texttt{lseu}), 4 \texttt{T1}
instances and 5 small classic instances. All have an integer leader and
integer follower variables. Each method has 60\,s per instance; an objective
equals the BOBILib optimum if it agrees to a relative $10^{-5}$, since BOBILib
prints six decimals.

\begin{table}[t]
\centering
\caption{Integer-follower methods on 80 BOBILib instances (SCIP 10.0, 60\,s).
``Feasible'': the method returned a bilevel-feasible objective (KKT answers
are checked with the follower MIP). ``$=$ opt'': equals the BOBILib optimum.}
\label{tab:rq3}
\begin{tabular}{lrrr}
\toprule
Method & Feasible & $=$ opt & Proved optimal \\
\midrule
KKT of the LP-relaxed follower & 65 & 16 & -- \\
Leader enumeration (Alg.~\ref{alg:enum}), warm-started & 79 & 49 & 35 \\
\bottomrule
\end{tabular}
\end{table}

\paragraph{RQ3.} Table~\ref{tab:rq3} summarizes the results. The KKT model
of the LP-relaxed follower returns a checked bilevel-feasible point on 65
instances, as Proposition~\ref{prop:relax} predicts, but only 16 of these
are optimal, and it reports infeasibility on 8 instances that have a bilevel
optimum. Both effects come from the same source: the model admits only
integer responses whose value equals the LP relaxation value. The leader
enumeration, started from the checked KKT value, proves optimality on 35
instances with a median of 2.5\,s and a maximum of 39.7\,s, and every proof
agrees with the BOBILib optimum. On the 44 instances where it stops at the
time limit its incumbent equals the optimum 14 times. No method returned an
objective below the BOBILib optimum, which is consistent with every
returned answer being bilevel feasible. One instance has an unbounded leader variable and is
outside the scope of the enumeration. The unproved instances are mostly the
larger DeNegre and Xu--Wang instances, where the HPR bound is weak and many
leader vectors must be visited; dedicated solvers such as
MibS~\cite{tahernejad2020branch} or the method of Fischetti et
al.~\cite{fischetti2017new} are the appropriate tools there.

\section{Discussion and Limitations}
\label{sec:limits}

The generated instances are small to medium and come from four families, so
the observed range of sufficient $M$ values is specific to them; the
propositions, not the sweep, are what carries over to other instances. The
follower check certifies bilevel feasibility, not optimality, and
Proposition~\ref{prop:anyM} shows it cannot reveal a too-small $M$ in the
KKT model; for that, a larger $M$ or an independent lower bound is needed.
The leader enumeration needs a pure-integer, bounded leader and its running
time grows with the number of leader vectors the HPR bound fails to exclude,
so it is a reference method for small and medium instances rather than a
replacement for dedicated bilevel solvers. All timings were taken on a
shared machine.

\section{Conclusion}
\label{sec:conclusion}

The big-M KKT model, stated with complementarity on both slacks and reduced
costs, always returns follower-optimal points; a wrong $M$ costs feasibility
or objective value, which we measured on 122 generated instances. For integer
followers, the KKT model of the LP relaxation supplies checked incumbents and
an HPR-bounded enumeration over leader decisions turns them into proven
optima on the BOBILib instances where it finishes. Both methods, the
instances and all result files are released with the paper.

\begin{thebibliography}{99}

\bibitem{balas1971intersection}
E.~Balas.
\newblock Intersection cuts---a new type of cutting planes for integer programming.
\newblock {\em Operations Research}, 19(1):19--39, 1971.

\bibitem{bolusani2024scip}
S.~Bolusani, M.~Besan{\c{c}}on, K.~Bestuzheva, et~al.
\newblock The {SCIP} Optimization Suite 9.0.
\newblock arXiv:2402.17702, 2024.

\bibitem{colson2007overview}
B.~Colson, P.~Marcotte, and G.~Savard.
\newblock An overview of bilevel optimization.
\newblock {\em Annals of Operations Research}, 153(1):235--256, 2007.

\bibitem{dempe2002foundations}
S.~Dempe.
\newblock {\em Foundations of Bilevel Programming}.
\newblock Kluwer Academic Publishers, 2002.

\bibitem{denegre2009branch}
S.~T. DeNegre and T.~K. Ralphs.
\newblock A branch-and-cut algorithm for integer bilevel linear programs.
\newblock In {\em Operations Research and Cyber-Infrastructure}, pages 65--78. Springer, 2009.

\bibitem{diasgarcia2024bileveljump}
J.~Dias~Garcia, G.~Bodin, and A.~Street.
\newblock {BilevelJuMP.jl}: Modeling and solving bilevel optimization problems in {Julia}.
\newblock {\em INFORMS Journal on Computing}, 36(2):327--335, 2024.

\bibitem{fischetti2016intersection}
M.~Fischetti, I.~Ljubi{\'c}, M.~Monaci, and M.~Sinnl.
\newblock Intersection cuts for bilevel optimization.
\newblock In {\em Integer Programming and Combinatorial Optimization (IPCO 2016)},
  LNCS 9682, pages 77--88. Springer, 2016.

\bibitem{fischetti2017new}
M.~Fischetti, I.~Ljubi{\'c}, M.~Monaci, and M.~Sinnl.
\newblock A new general-purpose algorithm for mixed-integer bilevel linear programs.
\newblock {\em Operations Research}, 65(6):1615--1637, 2017.

\bibitem{fischetti2018use}
M.~Fischetti, I.~Ljubi{\'c}, M.~Monaci, and M.~Sinnl.
\newblock On the use of intersection cuts for bilevel optimization.
\newblock {\em Mathematical Programming}, 172(1--2):77--103, 2018.

\bibitem{fortuny1981representation}
J.~Fortuny-Amat and B.~McCarl.
\newblock A representation and economic interpretation of a two-level programming problem.
\newblock {\em Journal of the Operational Research Society}, 32(9):783--792, 1981.

\bibitem{huangfu2018highs}
Q.~Huangfu and J.~A.~J. Hall.
\newblock Parallelizing the dual revised simplex method.
\newblock {\em Mathematical Programming Computation}, 10(1):119--142, 2018.

\bibitem{kleinert2020nofreelunch}
T.~Kleinert, M.~Labb{\'e}, F.~Plein, and M.~Schmidt.
\newblock There's no free lunch: on the hardness of choosing a correct big-{M} in bilevel optimization.
\newblock {\em Operations Research}, 68(6):1716--1721, 2020.

\bibitem{kleinert2021survey}
T.~Kleinert, M.~Labb{\'e}, I.~Ljubi{\'c}, and M.~Schmidt.
\newblock A survey on mixed-integer programming techniques in bilevel optimization.
\newblock {\em EURO Journal on Computational Optimization}, 9:100007, 2021.

\bibitem{moore1990mixed}
J.~T. Moore and J.~F. Bard.
\newblock The mixed integer linear bilevel programming problem.
\newblock {\em Operations Research}, 38(5):911--921, 1990.

\bibitem{pineda2019solving}
S.~Pineda and J.~M. Morales.
\newblock Solving linear bilevel problems using big-{M}s: not all that glitters is gold.
\newblock {\em IEEE Transactions on Power Systems}, 34(3):2469--2471, 2019.

\bibitem{tahernejad2020branch}
S.~Tahernejad, T.~K. Ralphs, and S.~T. DeNegre.
\newblock A branch-and-cut algorithm for mixed integer bilevel linear optimization problems and its implementation.
\newblock {\em Mathematical Programming Computation}, 12(4):529--568, 2020.

\bibitem{thurauf2026bobilib}
J.~Th{\"u}rauf, T.~Kleinert, I.~Ljubi{\'c}, T.~Ralphs, and M.~Schmidt.
\newblock {BOBILib}: Bilevel optimization (benchmark) instance library.
\newblock {\em Mathematical Programming Computation}, 18(3):819--876, 2026.

\end{thebibliography}

\end{document}
